\subsubsection{Sobolev 空间与 Fourier 变换}\label{subsec:sobolev-fourier-characterization}

现在考察整个 $\mathbb{R}^m$ 上的 Sobolev 空间. 显然, 由 Theorem~\ref{thm:fourier-inner-product-preservation} 可知, Fourier 变换将 $L^2(\mathbb{R}^m)$ 映入自身; 事实上, 它是 $L^2(\mathbb{R}^m)$ 中的等距映射. 此外, $D^\alpha u$ 的 Fourier 变换为 $(i\bm{\xi})^\alpha\hat u$. 因而立即得到以下结果.

\begin{Theorem}
\label{thm:sobolev-fourier-weighted-homeomorphism}
Fourier 变换 $\mathcal{F}$ 是从 $H^k(\mathbb{R}^m)$ 到加权空间 $L_w^2(\mathbb{R}^m)$ 上的同胚 (参见 Definition~\ref{def:weighted-lp-two-space}), 其中 $w(\bm{\xi})=1+|\bm{\xi}|^{2k}$.
\end{Theorem}

用记号 $L_k^2$ 表示这一加权 $L^2$ 空间. 容易看出 $\mathcal{S}(\mathbb{R}^m)$ 在 $L_k^2(\mathbb{R}^m)$ 中稠密. 于是 Theorem~\ref{thm:sobolev-fourier-weighted-homeomorphism} 表明, $\mathcal{S}(\mathbb{R}^m)$, 因而还有 $\mathcal{D}(\mathbb{R}^m)$, 在 $H^k(\mathbb{R}^m)$ 中稠密.

\begin{Corollary}
\label{cor:full-space-test-functions-sobolev-density}
$\mathcal{D}(\mathbb{R}^m)$ 在 $H^k(\mathbb{R}^m)$ 中稠密.
\end{Corollary}

该定理的另一应用是定义分数阶 Sobolev 空间.

\begin{Definition}
\label{def:fractional-order-sobolev-fourier}
若 $\mathcal{F}[u]$ 属于权为 $w(\bm{\xi})=1+|\bm{\xi}|^{2s}$ 的加权 $L^2$ 空间 $L_w^2(\mathbb{R}^m)=:L_s^2(\mathbb{R}^m)$, 则说 $u\in H^s(\mathbb{R}^m)$, $s\in\mathbb{R}^+$.
\end{Definition}

分数阶 Sobolev 空间有一种内蕴刻画, 基本上是 Hölder 连续性的 $L^2$ 类似物. 可以证明, $H^s(\mathbb{R}^m)$ 上的一个等价内积由下式给出:
\begin{equation}\label{eq:fractional-sobolev-intrinsic-inner-product}
\begin{aligned}
(u,v)_s&=(u,v)_{[s]}\\
&\quad+\sum_{|\alpha|=[s]}\int_{\mathbb{R}^m}\int_{\mathbb{R}^m}
\dfrac{(D^\alpha\bar u(\bm{x})-D^\alpha\bar u(\bm{y}))(D^\alpha v(\bm{x})-D^\alpha v(\bm{y}))}{|\bm{x}-\bm{y}|^{m+2(s-[s])}}\,\mathrm{d}\bm{x}\,\mathrm{d}\bm{y}.
\end{aligned}
\end{equation}
这里 $[s]$ 表示 $\leq s$ 的最大整数, \eqref{eq:fractional-sobolev-intrinsic-inner-product} 中的积分照例应解释为分别逼近 $u$ 和 $v$ 的光滑函数序列所对应的极限. 显然, 通过 \eqref{eq:fractional-sobolev-intrinsic-inner-product} 给出的定义可以立即推广到任意 $\Omega$.
